% \section{Evaluation Methodology}
% \label{sec:evaluation}

% \subsection{Algorithms Under Evaluation}
% % Более подробно про эксперименты: что мы забывали с помощью анлернинга 
% % перечислить еще раз алгоритмы, что мы подавали куда: GA, GD, NPO, RMU
% TODO: How we do unlearning with PopQA
% \subsection{Metrics}
% % Как считались метрики: 
% % Unlearning: Roudge-L Recall + LLM Juj
% % LLM: MMLU (acc), HellaSwag (norm_acc)
% TODO: Unlearning + LLM metrics info
% \subsection{Experimental Protocol}
% % Training details, hyper‑parameters, early‑stopping criteria, hardware footprint
% TODO: Training details, hyper‑parameters, early‑stopping criteria, hardware footprint


% \subsection{Evaluation Methodology}
% \label{sec:evaluation}
 \subsection{Evaluation Metrics}
\label{sec:evaluation}

% \subsection{Metrics}
In our research, we measure two key aspects: the success of the unlearning process itself and the preservation of the model's general utility.

\paragraph{Unlearning Efficacy and Locality.}
To quantify how well a model forgets a fact, following \cite{lin2004rouge, geng2025comprehensive} we use \textbf{ROUGE-L Recall}. This metric measures the word-level overlap between the model-generated response and the ground-truth answer.
\begin{itemize}
    \item A \textbf{lower} ROUGE-L score on the forget sets ($F_p^{\text{rare}}, F_p^{\text{pop}}$) indicates more effective unlearning.
    \item A \textbf{higher} (or close to the FT model) ROUGE-L score on the retain and validation sets ($R_p, V_{\text{DS}}, V_{\text{DA}}$) indicates better preservation of adjacent knowledge and thus higher unlearning locality.
\end{itemize}

\paragraph{General Capability Preservation.}
To ensure the unlearning process does not catastrophically damage the model's core abilities, we evaluate it on two standard downstream benchmarks:
\begin{itemize}
    \item \textbf{MMLU:} We report standard accuracy (\texttt{acc}) on this massive multitask benchmark to assess the model's broad world knowledge and problem-solving skills.
    \item \textbf{HellaSwag:} We report normalized accuracy (\texttt{norm\_acc}) for this commonsense reasoning task. This metric normalizes for the length of the possible continuations, providing a more robust measure of performance than standard accuracy.
\end{itemize}


\subsection{Unlearning Algorithms} %\& Evaluation
\label{sec:train}

Our benchmark contains a set of popular unlearning algorithms, which address the unlearning objective by optimizing the model's parameters $\theta$ to remove the association between a question $q$ and its specific answer $a$. For any given unlearning task, the target samples belong to a forget set $\mathcal{D}_f \in \{F_p^{\text{rare}}, F_p^{\text{pop}}\}$.

%\subsection{Unlearning Algorithms}\label{ssec:algorithms}
\textbf{Gradient Ascent (GA)} \cite{jang2022knowledge}. 
A foundational technique that reverses the learning process by maximizing the Negative Log-Likelihood (NLL) on the forget set. This is equivalent to performing gradient descent on the negative of the standard training loss. We define the NLL loss over a benchmark $\mathcal{D}$ as $\mathcal{L}_{\text{NLL}}(\mathcal{D}, \theta) = \frac{1}{|\mathcal{D}|} \sum_{(q,a) \in \mathcal{D}} (-\log P_{\theta}(a|q))$. The final loss for GA is the negative of this loss on the forget set:
\begin{equation}
\mathcal{L}^{\text{GA}}(\theta) = - \mathcal{L}_{\text{NLL}}(\mathcal{D}_f, \theta)
\end{equation}

\textbf{Gradient Difference (GD)} \cite{liu2022continual}.
This method extends GA by adding a regularization term to maintain performance on the retain set, $\mathcal{D}_r$. It simultaneously maximizes the loss on the forget data while minimizing it on the retain data, anchoring the model to its useful knowledge. The loss function is:
\begin{equation}
\mathcal{L}^{\text{GD}}(\theta) = -\mathcal{L}_{\text{NLL}}(\mathcal{D}_f, \theta) + \lambda \cdot \mathcal{L}_{\text{NLL}}(\mathcal{D}_r, \theta)
\end{equation}

\textbf{Negative Preference Optimization (NPO)} \cite{zhangnegative}.
NPO treats unlearning as a preference task. The objective is to make the probability of the undesired answer, $P_{\theta}(a|q)$, as small as possible relative to a reference model $P_{\theta_{\text{ref}}}(a|q)$. The loss function is formulated as:
\begin{equation}
\resizebox{0.85\linewidth}{!}{$
  \mathcal{L}^{\mathrm{NPO}}(\theta)
    = \frac{2}{\beta}\,
      \mathbb{E}_{(q,a)\in\mathcal{D}_f}
      \Bigl[\log\bigl(1+(\tfrac{P_{\theta}(a\mid q)}{P_{\theta_{\mathrm{ref}}}(a\mid q)})^{\beta}\bigr)\Bigr]
$}
\end{equation}

where $\beta$ is a hyperparameter, which we set to 1 as suggested in the original work.

\textbf{Representation Misdirection for Unlearning (RMU)} \cite{li2024wmdp}.
RMU operates on the model's internal hidden representations rather than its output probabilities. The core idea is to push the hidden states of the forget data towards a misdirection vector while ensuring the hidden states for the retain data remain close to those of the original model. The total loss is a weighted sum of a forget loss $\mathcal{L}_F$ and a retain loss $\mathcal{L}_R$.
\begin{gather}
\mathcal{L}_F = \mathbb{E}_{(q,a) \in \mathcal{D}_f}
\Biggl[ \frac{1}{|a|} \sum_{t \in a} \bigl\| h_{\theta}(t) - c \cdot u \bigr\|_2^2 \Biggr] \\
\mathcal{L}_R = \mathbb{E}_{(q,a) \in \mathcal{D}_r}
\Biggl[ \frac{1}{|a|} \sum_{t \in a} \bigl\| h_{\theta}(t) - h_{\theta_{\text{ref}}}(t) \bigr\|_2^2 \Biggr] \\
\mathcal{L}^{\mathrm{RMU}}(\theta) = \mathcal{L}_F + \lambda \cdot \mathcal{L}_R
\end{gather}

Here, $h_{\theta}(t)$ is the hidden state for token $t$ of the answer in the current model, $h_{\theta_{\text{ref}}}(t)$ is the state in the original reference model, $u$ is typically a pre-defined random vector for misdirection, and $c$ is a scaling factor.



